% !TeX root = ../artigo.tex
% Substitua as orientações abaixo pelo conteúdo da pesquisa.
\section{Introdução}
\label{sec:artigo-introducao}
Delimite o tema, o problema, a justificativa e o objetivo do estudo. A redação
científica precisa relacionar a pergunta de pesquisa às evidências que serão
apresentadas, com fundamentação nas fontes efetivamente consultadas
\citeind{severino2017}. O recorte do artigo deve caber em uma exposição concisa
sem omitir o método nem as limitações do trabalho.

\section{Fundamentação e trabalhos relacionados}
Apresente os conceitos necessários e compare trabalhos relacionados ao problema.
Esta seção pode ser integrada à introdução quando isso favorecer a argumentação.
Evite transformar o artigo em uma sequência de resumos de autores.

\section{Método}
Descreva o delineamento, os materiais ou dados, os procedimentos de análise e os
cuidados éticos aplicáveis. Em artigo tecnológico, explique também o projeto e a
avaliação do artefato. Informe condições suficientes para compreender como as
evidências foram produzidas.

\subsection{Procedimentos de avaliação}
Defina os critérios, as métricas e as condições da avaliação antes de discutir os
resultados. Diferencie evidência observada de resultado esperado. O exemplo
matemático a seguir demonstra a numeração independente de capítulos:
\begin{equation}
  p = \frac{n_{\mathrm{observado}}}{n_{\mathrm{total}}}.
  \label{eq:artigo-proporcao}
\end{equation}
Defina cada grandeza e explique sua pertinência; remova esta demonstração quando
não houver cálculo correspondente no estudo.

\section{Resultados e discussão}
Apresente os achados, relacione-os ao objetivo e confronte-os com a literatura.
Resultados e discussão podem formar seções distintas. Discuta limitações e
ameaças à validade. Figuras e tabelas devem apoiar a análise e ter fonte indicada.

\section{Considerações finais}
Retome a resposta ao problema, a contribuição e os limites do estudo. Indique
desdobramentos fundamentados nos resultados, sem introduzir dados inéditos nesta
seção. Este arquivo contém orientações de redação e não resultados de uma pesquisa.
